\documentclass[
	spanish,
	9pt,
	xcolor=table,
	handout,
	aspectratio=1610,
	ignorenonframetext
]{beamer}

\usepackage[spanish]{babel}
\usepackage[
	cursointroductorio,
	fullpagenumbering
]{../dunestyle-beamer}
\usepackage{booktabs}
\usepackage{longtable}
\usepackage{minted}

\graphicspath{{../images/}}

\usepackage[backend=biber,style=numeric, defernumbers=true, sorting=ynt,maxbibnames=4,maxcitenames=4]{biblatex}
\addbibresource{cmake-references.bib}

\title{
	Introducción a la herramienta cmake~\cmake{}
}
\subtitle{
	Parte I:
	\lstinline{bash},
	\lstinline{g++},
	\lstinline{cmake},
	\lstinline{gnuplot},
	\lstinline{gmsh},
	\lstinline{paraview}
}

\author{}
\institute[]
{
	\noindent
	Practique los ejemplos en Gitpod

	\href{https://gitpod.io\#/https://github.com/cpp-review-dune/dune-basics}{\includegraphics[width=4.5cm]{open-in-gitpod}}

	\vfill

	\begin{minipage}{0.12\paperwidth}
		Disponible en
	\end{minipage}%
	\hfill%
	\begin{minipage}{0.18\paperwidth}
		\href{https://github.com/cpp-review-dune}{\includegraphics[width=1.2cm]{GitHub-Mark-120px-plus}}
	\end{minipage}

	\vfill

	\begin{minipage}{0.26\paperwidth}
		¡Únete al grupo en Telegram!
	\end{minipage}%
	\hfill%
	\begin{minipage}{0.04\paperwidth}
		\href{https://t.me/joinchat/OsfYP1xnFlxjN2Ix}{\includegraphics[width=1.2cm]{telegram-logo}}
	\end{minipage}

}

\setbeamertemplate{bibliography item}{%
	\ifboolexpr{ test {\ifentrytype{book}} or test {\ifentrytype{mvbook}}
		or test {\ifentrytype{collection}} or test {\ifentrytype{mvcollection}}
		or test {\ifentrytype{reference}} or test {\ifentrytype{mvreference}} }
	{\setbeamertemplate{bibliography item}[book]}
	{\ifentrytype{online}
		{\setbeamertemplate{bibliography item}[online]}
		{\setbeamertemplate{bibliography item}[article]}}%
	\usebeamertemplate{bibliography item}}

\defbibenvironment{bibliography}
{\list{}
	{\settowidth{\labelwidth}{\usebeamertemplate{bibliography item}}%
		\setlength{\leftmargin}{\labelwidth}%
		\setlength{\labelsep}{\biblabelsep}%
		\addtolength{\leftmargin}{\labelsep}%
		\setlength{\itemsep}{\bibitemsep}%
		\setlength{\parsep}{\bibparsep}}}
{\endlist}
{\item}

\newcommand{\cmake}{%
	\begingroup\normalfont
	\includegraphics[height=2\fontcharht\font`B]{cmake}%
	\endgroup
}

\begin{document}

\frame[plain,noframenumbering,fragile]{
	\titlepage
}

\begin{frame}[fragile]
	\frametitle{CMake}
	Es un \emph{generador de sistemas de compilación} de código
	abierto e independiente del compilador o plataforma, es decir,
	produce instrucciones para otros sistemas de compilación como
	Makefile, Ninja, Visual Studio, Qt Creator, Android Studio y
	Xcode.
	También incluye características que permiten la instalación,
	empaquetamiento y soporte nativo de pruebas de software.

	\
	
	Resultados del aprendizaje:

	\begin{itemize}
		\item

		Construir diversos ejemplos de proyectos de CMake en los
		lenguajes C/C++/Fortran/Python que construyan ejecutables u
		objetos compartidos/estáticos/interfaces.

		\item

		Correr pruebas con \lstinline{ctest}, \lstinline{catch2},
		\lstinline{gtest} y \lstinline{pytest}.

		\item

		Usar dependencias de terceros en un proyecto CMake.
	\end{itemize}

	\

	Está orientado a cualquier estudiante, profesor, investigador,
	programador que quiera aprender a usar efectivamente CMake en un
	proyecto basado en Dune sobre Arch Linux.

	\

	No se asume que tenga experiencia en GNU/Linux o C++, pero
	sí que esté familiarizado con algún lenguaje de programación.

	% \lstinline{cmake} es una herramienta para automatizar la
	% generación de \lstinline{Makefile}.
	% TODO: Teoría
	% Explicar la necesidad de cmake (ver Scott, Rafal)
	% Instalación de cmake en Arch Linux.
	% Mencionar el manual pages de cmake
	% conceptos: build tree, source tree
	% Etapas de cmake con un diagrama.
	% TODO: Práctica
	% Ejemplo: un ejemplo de álgebra lineal con BLAS del libro de Simon.
	% Ejemplo: documentación con doxygen (Berner, Rafal).
	% Ejemplo con ctest: https://coderefinery.github.io/cmake-workshop/testing
\end{frame}

\begin{frame}[fragile]
	Tres herramientas de línea de comando:

	\begin{itemize}
		\item
		
		\href{https://man.archlinux.org/man/cmake.1}{\lstinline{/usr/bin/cmake}}

		\item
		
		\href{https://man.archlinux.org/man/cpack.1}{\lstinline{/usr/bin/cpack}}

		\item
		
		\href{https://man.archlinux.org/man/ctest.1}{\lstinline{/usr/bin/ctest}}
	\end{itemize}

	\

	Tres herramientas interactivas:

	\begin{itemize}
		\item
		
		\href{https://man.archlinux.org/man/ccmake.1}{\lstinline{/usr/bin/ccmake}}

		\item
		
		\href{https://man.archlinux.org/man/cmake-gui.1}{\lstinline{/usr/bin/cmake-gui}}
	\end{itemize}
\end{frame}

\begin{frame}[fragile]
	\frametitle{Instalación}
	Esta tecnología imprescindible se encuentra disponible en la
	mayoría de repositorios de distribuciones GNU/Linux importantes.
	Si desea instalar una versión actual, lo puede encontrar en el
	repositorio \href{https://archlinux.org/packages/extra/x86_64/cmake}{\lstinline{[extra]}}

	\begin{verbatim}
		[user@host somedir]$ sudo pacman -Syu
		[user@host somedir]$ sudo pacman -S cmake graphviz plantuml gcovr cppcheck python-cpplint ccache
	\end{verbatim}
\end{frame}

\begin{frame}[fragile]
	\frametitle{Un archivo \lstinline{CMakeLists.txt} minimal}
	Supongamos que tenemos el siguiente script \lstinline{CMakeLists.txt}

	\inputminted{cmake}{CMakeList.txt.sample}

	Entendamos, este script en más detalle.
	En este script definimos los requisitos para la construcción, desde
	el código fuente y objetivos, pasando por las pruebas, empaquetamiento, etc.
	Y delegaremos la tarea de compilación al programa \lstinline{make}.

	\begin{itemize}
		\item

		\href{https://cmake.org/cmake/help/latest/command/cmake_minimum_required.html}{\lstinline{cmake_minimum_required()}}
		indica la versión mínima de \lstinline{cmake} que requiere para ejecutar.

		\item
		
		\href{https://cmake.org/cmake/help/latest/command/project.html}{\lstinline{project()}}
		define el nombre del proyecto, su número de versión y que
		está escrito en el lenguaje de programación C++.
		
		\item

		\href{https://cmake.org/cmake/help/latest/command/set.html}{\lstinline{set()}}
		asigna la variable de entorno, en este caso establece la versión del estándar C++ 20.

		\item

		\href{https://cmake.org/cmake/help/latest/command/add_executable.html}{\lstinline{add_executable()}}
		crea un ejecutable a partir de un script en C++.
	\end{itemize}
\end{frame}

\begin{frame}[fragile]
	\frametitle{El proceso de construcción CMake}
	.
\end{frame}

\begin{frame}[fragile]
	\frametitle{Tipos de construcción CMake}

	\begin{itemize}
		\item
		
		Debug

		\item
		
		Release

		\item
		
		RelWithDebInfo

		\item
		
		MinSizeRel

		\item
		
		None
	\end{itemize}
\end{frame}

\begin{frame}[fragile]
	\frametitle{Generación automática de documentación con CMake}
	La documentación es esencial en cualquier proyecto exitoso.

	Veamos cómo integrar Doxygen con CMake para generar
	automáticamante la documentación para proyectos con CMake.

	CMake generará un archivo \lstinline{Doxyfile}.

	Esperamos que Doxygen genere la documentación de la interfaz de
	programación de aplicaciones (API) para cada clase y sus
	diagramas de herencia con \lstinline{dot} de graphviz.
\end{frame}

\begin{frame}[fragile]
	\frametitle{Pruebas con ctest}
	.
\end{frame}

% Reproducible builds, superbuilds

% TODO:
% Se va a separar una diapositiva de dune-cmake cuando esté terminado.
% Ejemplo: que sea de dune con dune-common, dune-geometry cmake.

\begin{frame}[fragile]
	\frametitle{El comando \href{https://gitlab.dune-project.org/core/dune-common/-/raw/master/bin/duneproject}{\lstinline{duneproject}}}

	Es un script asistente escrito en el lenguaje \lstinline{bash}
	que se encuentra en \lstinline{/usr/bin/duneproject}
	dentro del paquete \lstinline{dune-common}.
\end{frame}

\begin{frame}
	.
\end{frame}

\begin{frame}
	.
\end{frame}

\begin{frame}\transblindsvertical
	\frametitle{Referencias}
	%------------------------------------------------------------ 1
	\only<1>{
		\begin{itemize}
			\item Libros
			      \nocite{*}
			      \printbibliography[heading=none,keyword=book]
		\end{itemize}
	}
	%------------------------------------------------------------ 2
	\only<2>{
		\begin{itemize}
			\item Artículos
			      \printbibliography[heading=none,keyword=paper]
		\end{itemize}
	}
	%------------------------------------------------------------ 3
	\only<3>{
		\begin{itemize}
			\item Sitios web
			      \printbibliography[heading=none,keyword=online]
		\end{itemize}
	}
\end{frame}

\end{document}
%https://cliutils.gitlab.io/modern-cmake
%https://www.dune-project.org/sphinx/content/sphinx/core/index.html
%https://www.dune-project.org/sphinx/content/sphinx/core-2.7
%https://github.com/toeb/moderncmake/blob/master/Modern%20CMake.pdf